\documentclass[11pt,a4paper]{article}
% ============================================================
% Shared preamble for Claude Experiments reports (Ukrainian, LuaLaTeX)
% Compile with: lualatex report_*_uk.tex
% Include this AFTER \documentclass in each report:
%   \documentclass[11pt,a4paper]{article}
%   % ============================================================
% Shared preamble for Claude Experiments reports (Ukrainian, LuaLaTeX)
% Compile with: lualatex report_*_uk.tex
% Include this AFTER \documentclass in each report:
%   \documentclass[11pt,a4paper]{article}
%   % ============================================================
% Shared preamble for Claude Experiments reports (Ukrainian, LuaLaTeX)
% Compile with: lualatex report_*_uk.tex
% Include this AFTER \documentclass in each report:
%   \documentclass[11pt,a4paper]{article}
%   \input{../preamble_uk}
% ============================================================
% Cyrillic via LuaLaTeX + fontspec (pdflatex has no T2A here).
\usepackage{fontspec}
\setmainfont{DejaVu Serif}
\setsansfont{DejaVu Sans}
\setmonofont{DejaVu Sans Mono}
\usepackage[a4paper,top=2.5cm,bottom=2.5cm,left=2.5cm,right=2.5cm]{geometry}
\usepackage{amsmath,amssymb}
\usepackage{graphicx}
\usepackage{subcaption}
\usepackage{booktabs}
\usepackage{xcolor}
\usepackage[colorlinks=true,allcolors=blue]{hyperref}
\usepackage{microtype}
\usepackage{float}
\usepackage{array}
\microtypesetup{expansion=false}

% Color scheme
\definecolor{colMicro}{HTML}{222222}
\definecolor{colExact}{HTML}{2878b5}
\definecolor{colPlacement}{HTML}{e07a1f}
\definecolor{colBoundary}{HTML}{429c74}
\definecolor{colCurrent}{HTML}{c23b55}
\definecolor{colMM}{HTML}{c23b55}

\providecommand{\SIG}{\textcolor{green!50!black}{\textbf{*}}}
\providecommand{\NS}{\textcolor{gray}{n.s.}}
\providecommand{\DRMS}{\Delta\mathrm{RMST}}
\providecommand{\ci}[2]{$[{#1},\;{#2}]$}

% ============================================================
% Cyrillic via LuaLaTeX + fontspec (pdflatex has no T2A here).
\usepackage{fontspec}
\setmainfont{DejaVu Serif}
\setsansfont{DejaVu Sans}
\setmonofont{DejaVu Sans Mono}
\usepackage[a4paper,top=2.5cm,bottom=2.5cm,left=2.5cm,right=2.5cm]{geometry}
\usepackage{amsmath,amssymb}
\usepackage{graphicx}
\usepackage{subcaption}
\usepackage{booktabs}
\usepackage{xcolor}
\usepackage[colorlinks=true,allcolors=blue]{hyperref}
\usepackage{microtype}
\usepackage{float}
\usepackage{array}
\microtypesetup{expansion=false}

% Color scheme
\definecolor{colMicro}{HTML}{222222}
\definecolor{colExact}{HTML}{2878b5}
\definecolor{colPlacement}{HTML}{e07a1f}
\definecolor{colBoundary}{HTML}{429c74}
\definecolor{colCurrent}{HTML}{c23b55}
\definecolor{colMM}{HTML}{c23b55}

\providecommand{\SIG}{\textcolor{green!50!black}{\textbf{*}}}
\providecommand{\NS}{\textcolor{gray}{n.s.}}
\providecommand{\DRMS}{\Delta\mathrm{RMST}}
\providecommand{\ci}[2]{$[{#1},\;{#2}]$}

% ============================================================
% Cyrillic via LuaLaTeX + fontspec (pdflatex has no T2A here).
\usepackage{fontspec}
\setmainfont{DejaVu Serif}
\setsansfont{DejaVu Sans}
\setmonofont{DejaVu Sans Mono}
\usepackage[a4paper,top=2.5cm,bottom=2.5cm,left=2.5cm,right=2.5cm]{geometry}
\usepackage{amsmath,amssymb}
\usepackage{graphicx}
\usepackage{subcaption}
\usepackage{booktabs}
\usepackage{xcolor}
\usepackage[colorlinks=true,allcolors=blue]{hyperref}
\usepackage{microtype}
\usepackage{float}
\usepackage{array}
\microtypesetup{expansion=false}

% Color scheme
\definecolor{colMicro}{HTML}{222222}
\definecolor{colExact}{HTML}{2878b5}
\definecolor{colPlacement}{HTML}{e07a1f}
\definecolor{colBoundary}{HTML}{429c74}
\definecolor{colCurrent}{HTML}{c23b55}
\definecolor{colMM}{HTML}{c23b55}

\providecommand{\SIG}{\textcolor{green!50!black}{\textbf{*}}}
\providecommand{\NS}{\textcolor{gray}{n.s.}}
\providecommand{\DRMS}{\Delta\mathrm{RMST}}
\providecommand{\ci}[2]{$[{#1},\;{#2}]$}


\title{CE-03: Похибка MicroMacro залежно від Beta Ratio\\
       \large Звіт про експеримент --- проєкт Two-Scale CTMC}
\author{Claude Experiments Pipeline}
\date{2026-07-31}

\begin{document}
\maketitle

\begin{abstract}
Цей звіт кількісно оцінює, як відношення $r = \beta_{\text{out}}/\beta_{\text{in}}$
між інтер- та інтракомунітарними інтенсивностями інфікування впливає на похибку
апроксимації MicroMacro. П'ять сценаріїв проходять по $r \in \{0.1, 0.3, 1.0, 2.0, 5.0\}$
на фіксованій топології chain-of-cliques ($L=10$, $n=100$, $d=6$, $k=5$)
зі 100 парними реплікаціями в кожному.
Головний результат --- \textbf{зміна знаку} у $\DRMS$:
MicroMacro прибуває \emph{пізніше} за повний Gillespie при малому $r$
($\DRMS = {+}4796$ при $r=0.1$) і \emph{раніше} при великому $r$
($\DRMS = -1355$ при $r=1.0$, $-1638$ при $r=2.0$).
Crossover відбувається між $r=0.3$ та $r=1.0$.
При $r \ge 2$ епідемія в повному Micro рідко досягає останньої громади
в межах горизонту симуляції (completion 54\% та 1\% відповідно),
тоді як MicroMacro завершується в усіх прогонах --- що вказує на якісний
false-positive bias за сильного інтеркомунітарного зв'язку.
\textit{Ukrainian summary}: див.\ \texttt{docs/ce03\_beta\_ratio/report.md}.
\end{abstract}

\section{Дизайн експерименту}

\subsection{Мета}
Визначити, як відношення інтеркомунітарної до інтракомунітарної інтенсивності
інфікування модулює похибку апроксимації MicroMacro, і охарактеризувати напрям
(передчасне vs.\ запізніле прибуття) на всьому діапазоні параметра.

\subsection{Мережа та параметри}
\begin{itemize}
  \item Топологія: chain-of-cliques, $L=10$ громад, по $n=100$ вузлів у кожній,
        степінь $d=6$, $k=5$ інтеркомунітарних ребер на межу.
  \item Інтракомунітарна інтенсивність інфікування: $\beta_{\text{in}} = 0.003$.
  \item Інтеркомунітарна інтенсивність інфікування: $\beta_{\text{out}} = r \cdot \beta_{\text{in}}$.
  \item Відношення інтенсивностей: $r \in \{0.1,\; 0.3,\; 1.0,\; 2.0,\; 5.0\}$.
  \item Одужання: $\gamma = 0$ (чиста модель SI).
  \item Реплікації: 100 парних прогонів на сценарій (та сама мережа й seed).
  \item Горизонт $T_{\text{end}}$: залежить від сценарію (20{,}000 до 2{,}000;
        коротший для більшого $r$, щоб час обчислень залишався прийнятним).
\end{itemize}

\subsection{Первинна метрика}
\[
  \DRMS = \overline{T_{\text{MM}}} - \overline{T_{\text{Micro}}}
\]
де часи прибуття використовують restricted-mean оцінки з censoring на
$T_{\text{end}}$ (цензуровані прогони дають внесок $T_{\text{end}}$ у середнє).
Додатний $\DRMS$ означає, що MicroMacro прогнозує \emph{запізніле} прибуття;
від'ємний --- \emph{передчасне} прибуття.
Bootstrap 95\% довірчі інтервали використовують 10{,}000 ресемплів.
Відносна відстань Wasserstein $\text{rel-}W_1$ вимірює розбіжність розподілів,
нормовану на середній час прибуття Micro.

\section{Результати}

\begin{table}[H]
\centering
\caption{Результати CE-03: похибка MicroMacro за beta ratio.
  $\DRMS$ = різниця цензурованих середніх (MicroMacro $-$ Micro).
  $r = \beta_{\text{out}}/\beta_{\text{in}}$.
  ``Comp.'' = частка прогонів Micro, що досягли останньої громади
  до $T_{\text{end}}$ (MicroMacro: 1.00 для всіх сценаріїв).
  Усі CI виключають нуль.}
\label{tab:ce03-main}
\begin{tabular}{rrrrrrc}
\toprule
$r$ & $\beta_{\text{out}}$ & $\DRMS$ & \multicolumn{2}{c}{95\% CI} & rel-$W_1$ & Micro comp. \\
\cmidrule(lr){4-5}
 & & & lower & upper & & \\
\midrule
0.1 & 0.0003 & $+4796$ & $+4384$ & $+5213$ & 1.624 & 1.00 \\
0.3 & 0.0009 & $+562$  & $+300$  & $+812$  & 0.297 & 1.00 \\
1.0 & 0.0030 & $-1355$ & $-1494$ & $-1220$ & 0.456 & 1.00 \\
2.0 & 0.0060 & $-1638$ & $-1746$ & $-1527$ & 0.558 & 0.54 \\
5.0 & 0.0150 & $-1213$ & $-1283$ & $-1141$ & 0.566 & 0.01 \\
\bottomrule
\end{tabular}
\end{table}

\subsection{Ключові спостереження}

\begin{enumerate}
  \item \textbf{Зміна знаку.}  $\DRMS$ додатний при $r \le 0.3$ і
        від'ємний при $r \ge 1.0$.  Crossover $r^*$ лежить між 0.3 та 1.0.
  \item \textbf{Велика магнітуда на краях.}  $|\DRMS| \approx 4800$ при $r=0.1$
        і $\approx 1600$ при $r=2.0$ --- обидві великі відносно горизонту.
  \item \textbf{Розбіжність розподілів.}  $\text{rel-}W_1 = 1.624$ при $r=0.1$
        (розподіли відрізняються більше, ніж на середній час прибуття Micro).
  \item \textbf{Колапс completion при великому $r$.}  При $r=2.0$ лише 54\%
        прогонів повного Micro досягають останньої громади в межах $T_{\text{end}}=3000$;
        при $r=5.0$ лише 1\%.  MicroMacro завжди повідомляє про повне проходження,
        даючи false-positive рівень епідемії 99 процентних пунктів при $r=5.0$.
\end{enumerate}

\section{Інтерпретація}

\subsection{Механізм при малому $r$ (запізніле прибуття, $\DRMS > 0$)}

Коли $\beta_{\text{out}} \ll \beta_{\text{in}}$, інтракомунітарне поширення швидке,
а інтеркомунітарне --- рідкісне.  У повній моделі Gillespie (Micro)
$k=5$ gateway-вузлів на межі вихідної громади можуть спрацьовувати своїми
cross-community ребрами щойно \emph{вони самі} стають інфікованими,
незалежно від загальногромадської частки інфікованих $\bar{i}(t)$.

MicroMacro використовує агрегований import hazard
$\lambda_{\text{macro}}(t) = \beta_{\text{out}} \cdot \bar{i}(t)$,
пропорційний до загальногромадського середнього.  Коли $\beta_{\text{out}}$
крихітний, цей hazard залишається незначним, доки $\bar{i}(t)$ суттєво не зросте,
тож macro import event спрацьовує \emph{пізно}.  Прямі gateway-ребра в Micro
можуть спрацювати раніше, даючи $\DRMS > 0$.

\subsection{Механізм при великому $r$ (передчасне прибуття, $\DRMS < 0$)}

Коли $\beta_{\text{out}} \ge \beta_{\text{in}}$, агрегований hazard
$\lambda_{\text{macro}}(t) = \beta_{\text{out}} \cdot \bar{i}(t)$ стає
великим щойно в вихідній громаді з'являються будь-які інфіковані.  MicroMacro
агресивно спрацьовує imports і розміщує кожен у рівномірно випадковому вузлі
цільової громади, штучно прискорюючи проходження ланцюга.

У повному Micro епідемія має поширюватися лише через $k=5$ фізичних ребер
на межу.  Навіть за великого $\beta_{\text{out}}$ проходження дев'яти
послідовних меж у ланцюзі обмежене внутрішньогромадським накопиченням
на кожному кроці.  Micro рідко завершує ланцюг у межах короткого $T_{\text{end}}$
(1\% при $r=5.0$), тоді як MicroMacro завжди завершує.

\subsection{Порівняння з baseline CE-01}

CE-01 використовував однорідні інтенсивності ($r=1$) і виявив $\DRMS = -209$ із
rate abstraction як домінантним джерелом похибки.  CE-03 при $r=1.0$ дає
$\DRMS = -1355$; більша магнітуда узгоджується з коротшим
$T_{\text{end}} = 4000$ (проти 5000 в CE-01).  Обидва підтверджують передчасне
прибуття при $r=1$.

\section{Висновок}

CE-03 виявляє якісну зміну напряму похибки MicroMacro як функції
від $r$:

\begin{itemize}
  \item \textbf{Мале $r$ ($r < r^* \approx 0.5$--$0.8$):}
        MicroMacro прогнозує \emph{запізніле} прибуття.  Агрегований hazard
        недооцінює ранню активність gateway, через що macro import спрацьовує занадто пізно.
  \item \textbf{Велике $r$ ($r > r^*$):}
        MicroMacro прогнозує \emph{передчасне} прибуття.  Агрегований hazard
        переоцінює ефективну інтенсивність перетину багатьох меж.
        При $r=5.0$ якісний прогноз є хибним: Micro рідко поширюється
        через увесь ланцюг, тоді як MicroMacro завжди (99\% false-positive rate).
  \item \textbf{Crossover поблизу $r^* \approx 0.5$--$0.8$:}
        Обидва напрями похибки взаємно скорочуються.  Це скорочення залежить від
        топології, а не є режимом коректної апроксимації.
\end{itemize}

Одна скалярна корекція $\beta_{\text{out}}$ не може прибрати bias
на всьому діапазоні параметрів.  Потрібні структурні корекції --- як-от
використання фактичної кількості інфікованих gateway-вузлів замість
громадського середнього.

\section*{Розташування даних}
\texttt{output/claude\_experiments/ce03\_beta\_ratio/summary.csv}

\section*{Пов'язані експерименти}
CE-01 (ablation: rate abstraction vs.\ placement), CE-04 (SIR: ефект одужання).
\end{document}
